
Reinforcement learning from human feedback (RLHF) has become standard for aligning language models, yet evaluation metrics focus on AI-side properties (helpfulness, harmlessness) with limited attention to human-side effects such as agency preservation. This study tests whether linguistic markers validated in human psychology (modal verbs, hedging language, alternative-framing phrases) can serve as computational proxies for measuring agency preservation in RLHF contexts. Using 169,352 preference pairs from the Anthropic HH-RLHF dataset, we conducted a two-stage validation: (1) feasibility testing showed that markers can be reliably extracted with high precision (100\%) and sufficient variance (CV=0.781), but (2) construct validation revealed that markers do not validly measure agency preservation, failing effect size criteria (Cohen's d=-0.018, 88\% below threshold of 0.15), internal consistency (Cronbach's $\alpha =0.42$ vs. required 0.7), and cross-split replication (0/2 splits passed). These findings demonstrate that statistical significance (p<0.001) does not guarantee practical relevance with large samples (N=169K), and that cross-domain proxy transfer from human psychology to AI-generated text requires empirical validation rather than theoretical assumption. The negative result establishes validation requirements for future computational operationalization of bidirectional alignment metrics.
